%%%%%%%%%%%%%%%%%%%%%%%%%%%%%%%%%%%%%%%%%%%%%%%%%%%%%%%%%%%%%%%%%%%%%%%%%%%%%%%%%%%%%%%
%%%%%%%%%%%%%%%%%%%%%%%%%%%%%%%%%%%%%%%%%%%%%%%%%%%%%%%%%%%%%%%%%%%%%%%%%%%%%%%%%%%%%%%%
% Beamer Template - Version: Leuphana
% © Fynn Lohre & Leuphana for Graphics
%
% The following template is designed for a presentation in Leuphana style.
%
% In case of any questions - feel free to reach out: VARFynn@gmail.com
% If this template helps you, endorsements on LinkedIn or in persona are welcome.
%
%%%%%%%%%%%%%%%%%%%%%%%%%%%%%%%%%%%%%%%%%%%%%%%%%%%%%%%%%%%%%%%%%%%%%%%%%%%%%%%%%%%%%%%%

%---------------------------------------------------------------------------------------

%%%%%%%%%%%%%%%%%%%%%%%%%%%%%%%%%%%%%%%%%%%%%%%%%%%%%%%%%%%%%%%%%%%%%%%%%%%%%%%%%%%%%%%%
%
% 1.  Preamble
%
% This defines the properties of the document.
% Includes: Font Size, Spacing & other "Style Elements".
%
%%%%%%%%%%%%%%%%%%%%%%%%%%%%%%%%%%%%%%%%%%%%%%%%%%%%%%%%%%%%%%%%%%%%%%%%%%%%%%%%%%%%%%%%

%% 1.1 Basic Packages

\documentclass[11pt,aspectratio=169] {beamer}
\usetheme{CambridgeUS}
\usepackage[utf8]{inputenc}
\usepackage{babel}
\usepackage{amsmath}
\usepackage{amsfonts}
\usepackage{amssymb}
\usepackage{mathtools}
\usepackage{tikz}
\setbeamercovered{transparent} 
%\setbeamertemplate{navigation symbols}{} 
\setbeamertemplate{caption}[numbered]
\setbeamertemplate{theorems}[numbered] 


%% 1.2 Basic Packages Color Scheme Leuphana

\definecolor{leuphana-gray}{RGB}{100,100,100} % Define a darker gray color
\setbeamercolor{section in toc}{fg=black,bg=white}
\setbeamercolor{alerted text}{fg=black} % Set alerted text to black
\setbeamercolor*{palette primary}{fg=black,bg=leuphana-gray!15!white} % Use the darker gray for the background
\setbeamercolor*{palette secondary}{fg=black,bg=leuphana-gray!2!white}
\setbeamercolor*{palette tertiary}{bg=leuphana-gray!40!white} % Use the darker gray for the background
\setbeamercolor*{palette quaternary}{fg=black,bg=leuphana-gray!5!white}
\setbeamercolor*{sidebar}{fg=black,bg=leuphana-gray!15!white}
\setbeamercolor*{palette sidebar primary}{fg=leuphana-gray!10!black} % Use the darker gray for the foreground
\setbeamercolor*{palette sidebar secondary}{fg=black}
\setbeamercolor*{palette sidebar tertiary}{fg=leuphana-gray!50!black} % Use the darker gray for the foreground
\setbeamercolor*{palette sidebar quaternary}{fg=leuphana-gray!10!white}
\setbeamercolor{frametitle}{fg=black}
\setbeamercolor{title}{fg=black}


%% 1.3 Pre Info for Titlepage

\title[Short Title for Fooder]{Titel}
\author[Short Name for Fooder]{John Doe}
\institute[]{\includegraphics[scale=0.06]{Leuphana.svg.png}}

\addtobeamertemplate{frametitle}{}{%
\begin{tikzpicture}[remember picture,overlay]
  \node at (current page.north east) [anchor=north east, inner sep=16pt] {\includegraphics[scale=0.14]{Logo_Universität_Lüneburg.png}};
\end{tikzpicture}}


%---------------------------------------------------------------------------------------

%%%%%%%%%%%%%%%%%%%%%%%%%%%%%%%%%%%%%%%%%%%%%%%%%%%%%%%%%%%%%%%%%%%%%%%%%%%%%%%%%%%%%%%%
%
% 2. Slides
%
%%%%%%%%%%%%%%%%%%%%%%%%%%%%%%%%%%%%%%%%%%%%%%%%%%%%%%%%%%%%%%%%%%%%%%%%%%%%%%%%%%%%%%%%

\begin{document}

% Title Page
\begin{frame}
\titlepage
\end{frame}

% TOC 
\begin{frame}{Table of Contents}
\tableofcontents
\end{frame}
\section{Type Section I  Here}
\subsection{Type Subsection I Here}

\begin{frame}{Table of Contents II Version}
 \tableofcontents
    [
        currentsection,
        currentsubsection,
        subsectionstyle=show/shaded/hide
    ]
\end{frame}

% Equations slide
\begin{frame}{Equations}
\[
E=mc^2
\]
\begin{equation}
\hat{\beta} = 
\end{equation}
\end{frame}

% Double-Sided slide
\begin{frame}[allowframebreaks]{Double-Sided Slide}
This slide is split into multiple frames.
\framebreak
Second part of the slide.
\end{frame}

% Block slide
\begin{frame}{Block}
\begin{block}{Block Title}
This is a block.
\end{block}
\begin{alertblock}{Alert Block Title}
This is an alert block.
\end{alertblock}
\begin{exampleblock}{Example Block Title}
This is an example block.
\end{exampleblock}
\end{frame}

% Theorem slide
\begin{frame}
\begin{Theorem}
XYZ 
\end{Theorem}
\end{frame}

% Overlay slide
\begin{frame}{Overlay}
\begin{itemize}
\item<1-> Item 1
\item<2-> Item 2
\item<3-> Item 3
\end{itemize}
\end{frame}

% Image and Text slide
\begin{frame}{Image and Text}
\begin{columns}
\begin{column}{0.4 \textwidth}
\begin{figure}
\includegraphics[width=\textwidth]{example-image}
\caption{Caption}
\end{figure}
\end{column}
\begin{column}{0.4\textwidth}
Some text.
\end{column}
\end{columns}
\end{frame}

% Bullet Points slide
\begin{frame}{Bullet Points \& Enumerate}
\begin{itemize}
\item Bullet point 1
\item Bullet point 2
\item Bullet point 3
\end{itemize}
\begin{enumerate}
\item Bullet point 1
\item Bullet point 2
\item Bullet point 3
\end{enumerate}
\end{frame}

% Table slide
\begin{frame}{Table}
\begin{tabular}{|c|c|}
\hline
A & B \\
\hline
1 & 2 \\
3 & 4 \\
\hline
\end{tabular}
\end{frame}

% Top Centered Slide
\begin{frame}[t]{Top Centered Slide}
This slide is centered at the top.
\end{frame}

%E pochal Slide
\begin{frame}{Epochal Slide}
\centering
\textbf{\Huge Epochal Event}
\end{frame}

% Quote Slide
\begin{frame}{Quote Slide}
\begin{quote}
This is a quote.
\end{quote}
\end{frame}

% Code Slide
\begin{frame}[fragile]{Code Slide}
\begin{verbatim}
for i in range(5):
    print(i)
\end{verbatim}
\end{frame}
\section{Type Section II Here}

% Image Grid Slide
\begin{frame}
\frametitle{Image Grid Slide}
\begin{figure}
\includegraphics[width=0.3\textwidth]{example-image-a}
\includegraphics[width=0.3\textwidth]{example-image-b}
\includegraphics[width=0.3\textwidth]{example-image-c}
\caption{Images in a grid}
\end{figure}
\end{frame}

%Blank Frame
\bgroup
\setbeamercolor{background canvas}{bg=black}
\begin{frame}[plain]{}
\end{frame}
\egroup

\section{Backup}
% Two Columns Slide with Divider
\begin{frame}[t]
\frametitle{Two Columns Slide with Divider}
\begin{columns}[T]
\column{0.46\textwidth}
\underline{\textbf{Column 1 Subhead}}
\begin{itemize}
\item Text
\item Text
\end{itemize}

\column{0.02\textwidth}
\centering
\rule{0.5pt}{4cm} % Vertical divider

\column{0.46\textwidth}
\underline{\textbf{Column 2 Subhead}}
\begin{itemize}
\item Text
\item Text
\end{itemize}
\end{columns}
\vfill
\begin{block}{Takeaway Box}
Text
\end{block}
\vfill
\end{frame}

\end{document}